% Lösungen Einheit 4 von 6 – Symmetrie (zusammenbau v0.1)
\einheitenkopf{Einheit 4 von 6 · Symmetrie}

\erg{17}{a) achsensymmetrisch \quad b) achsensymmetrisch zur y-Achse: nur gerade Exponenten, das Absolutglied zählt gerade \quad c) $f(-x) = (-x)^2 \cdot e^{-(-x)^2} = x^2 \cdot e^{-x^2} = f(x)$: achsensymmetrisch zur y-Achse \quad d) $f_0(x) = -x$: nur ungerader Exponent, punktsymmetrisch zum Ursprung \quad e) $S_y(0 | 7)$; achsensymmetrisch zur y-Achse, nur gerade Exponenten \quad f) Nullstelle $-2{,}5$; Hochpunkt $H'(-1{,}5 | 4)$ \quad g) Fußpunkte $(-12 | 0)$ und $(12 | 0)$, höchster Punkt $(0 | 6)$ \quad h) Schnittstellen $-3$ und $3$; $c = f(3) = 4{,}5$ \quad i) achsensymmetrisch zur y-Achse: nur gerade Exponenten, das Absolutglied zählt gerade; Grad gerade, Leitkoeffizient positiv: für $x \to -\infty$ und für $x \to +\infty$ gilt $f(x) \to +\infty$ – der negative Koeffizient bei $x^2$ spielt dafür keine Rolle}
\erg{18}{a) Fehler: das Absolutglied zählt als gerader Exponent (x hoch null). Richtig: nur gerade Exponenten, achsensymmetrisch zur y-Achse \quad b) Das Absolutglied ist die Zahl mal x hoch null, und null ist gerade; setzt man $-x$ ein, ändert es sich nicht.}
